\section{Claims}
\label{sec:claims}

\begin{description}[leftmargin=2.2em,style=nextline]
  \item[\textbf{C1 (H2, headline).}]
    Because the model \emph{is} the prior, model error enters the analysis
    unbuffered, and its cost is not a constant offset: it \emph{collapses the
    marginal value of observations} ($6.2\times$ for strong-constraint 4D-Var).
    This is the mechanism behind DA's limited return on sparse observations.
    \needsrun{C1 is argued against strong-constraint and filtering methods only
    until weak-constraint 4D-Var is benchmarked (experiment D0). Weak-constraint
    formulations explicitly relax H2 and are the classical methods \emph{allowed}
    to know about model error. If D0 shows weak-constraint 4D-Var closes most of
    the gap, C1 must be narrowed accordingly --- which is itself a publishable
    result.}
  \item[\textbf{C2 (H3a).}]
    The ODE state representation fixes the conditioning of the inverse problem.
    Identical physics in a different state variable produces large, predictable
    skill differences via $q \approx \nabla^2\psi$ and its $K^2$ error
    amplification.
    \todo{Decide whether the QG result carries C2 alone or whether the Lorenz-96
    slow/fast analogue (D4) is needed first.}
  \item[\textbf{C3 (H1).}]
    Markovian factorisation forecloses adaptation to the observation
    configuration. A configuration-distribution-aware estimator dominates both
    configuration-specific and configuration-agnostic ones (0.473 / 1.081 at
    full / zero fast-variable density, against 0.486 / 1.292 and 0.537 / 1.082).
  \item[\textbf{C4 (H1 $+$ H3a $\Rightarrow$ S-b).}]
    The ensemble-Kalman class carries $\PsiNG \equiv 0$ by construction, so its
    posterior characterisation has a \emph{structural} ceiling --- and the
    alternatives measured so far buy calibration or accuracy, never both.
  \item[\textbf{C5 (attribution).}]
    Each hypothesis can be relaxed separately and the resulting gain attributed
    to a named term of the operator partition. This is what makes the paper a
    diagnosis rather than a benchmark.
  \item[\textbf{C6 (H2 $\times$ H3b).}]
    The ODE representation's sensitivity to $(\theta, z)$ is \emph{not a proxy}
    for the sensitivity of $p(\xv \mid \yv, \theta, z)$ to $(\theta, z)$. DA
    conflates them because its estimator is the model; a conditioned amortised
    estimator separates them; and the gap is large --- $1.68$--$1.80\times$
    against $\leq 1.5\%$.
  \item[\textbf{C7 (H4, classical corollary).}]
    The value of the adjoint is \emph{conditional on model fidelity}. At \Sz{}
    the adjoint method and the best ensemble filter convert extra observations
    into skill at the same rate; under model error the adjoint method's return
    collapses $3.3\times$ harder.
\end{description}

\section{Remaining experiments}
\label{sec:experiments}

The claims above rest on evidence already in hand. The following are designed
and, in most cases, implemented but not yet run. They are listed because the
paper's final form depends on them.

\begin{description}[leftmargin=2.2em,style=nextline]
  \item[\textbf{D0 --- weak-constraint 4D-Var on Lorenz-96.}] \emph{Prerequisite;
    blocks C1.} The implementation exists, including a whitened-control
    formulation with gradient sanitisation and NaN recovery. The investigation
    behind it established that the naive formulation's failure on this system is
    an \emph{optimiser} problem, not a property of weak-constraint DA: Adam
    leaves the observation cost two to three orders of magnitude above its noise
    floor through a 500-step chaotic unroll, whereas L-BFGS with the
    strong-constraint method's own validated settings reaches it.
    \needsrun{The \Sz{}/\So{} benchmark run itself. Highest value, lowest cost
    item in this plan.}
  \item[\textbf{D1 --- H1 isolated.}] Filter vs.\ window-smoother vs.\
    window-amortised at matched model, observations and compute.
  \item[\textbf{D2 --- H2 dose--response on a common axis.}] Sweep model error
    continuously rather than the binary \Sz{}/\So{}, giving the \emph{same}
    biased parameter set to the classical method's forward model and to the
    learned scheme's conditioning inputs.
  \item[\textbf{D3 --- the marginal-value surface.}]
    $\mathrm{d}\,\text{skill}/\mathrm{d}\,\text{density}$ over a two-dimensional
    grid of observation density $\times$ model error, for every method class on
    matched axes. Table~\ref{tab:marginal} is one row of this surface.
    \caveat{The classical density study varies \emph{temporal} density and the
    learned one varies \emph{channel} density. These are different axes and must
    not be tabled together until D3 puts them on one protocol.}
  \item[\textbf{D4 --- H3a by re-parameterisation.}] The QG $\psi/q$ result
    generalised, relating the skill difference to the conditioning of the map
    between representations.
  \item[\textbf{D5 --- S-b done properly.}] Rank histograms, spread/skill against
    $\sqrt{(N-1)/(N+1)}$, and CRPS/ES at a single uniform ensemble convention
    across all method classes. The largest new-code item.
  \item[\textbf{D6 --- attribution.}] Report $\Psimean$, $\PsiG$, $\PsiNG$ per
    method per regime, on both axes and in both point-estimate and dispersion
    terms.
  \item[\textbf{D7 --- the two sensitivities on one grid.}] See
    Section~\ref{sec:sensitivity}.
\end{description}
